% Siden den industrielle revolution har fossile brændstoffer været menneskets primære energikilde.
% Forbrænding af fossile brændstoffer frigiver store mængder drivhusgasser, som bidrager til global opvarmning og klimaforandringer. 
% Derfor er der stigende fokus på alternative energikilder som sol, vind og kernekraft.

% Kernekraft udnyttes i dag primært gennem fission, hvor tunge atomkerner sammen spaltes og frigiver energi. 
% Fission har det negative biprodukt, radioaktivt affald, som medfører en række sikkerheds/miljømæssige udfordringer.

% Fusion, hvor lette atomkerner sammensmeltes, betragtes som en lovende fremtidig energikilde. 
% Processen kan potentielt levere store mængder energi hvorved affaldet har meget højere halveringstid.

% For at opnå fusion opvarmes en kilde af atomer til ekstremt høj temperatur, hvor stoffet overgår 
% til plasma, en ioniseret gas bestående af frie elektroner og ioner. 
% I moderne fusionsanlæg som tokamaker fastholdes plasmaet ved hjælp af stærke elektro magnetiskefelter,
% så det ikke kommer i kontakt med tokamakkens vægge, hvilket medfører energitab og ustabilitet.

En central udfordring inden for fusionsforskning er at forstå og beskrive plasmaets dynamik, 
særligt i randzonen, hvor turbulens og transportfænomener har stor betydning for energitab og stabilitet.
Disse processer er moddeleret med koblede partielle differentialligninger (PDE'er). 
Projektet fokuserer på hot edge (scrape of layer) electrostatic (HESEL) modellering. 


Til at simulere HESEL-moddelen bruges numerisk finite difference. Dette er et centralt værktøj i moderne plasmafysik.
Ulæmpen ved denne tilgang er at simuleringerne er beregningstunge og skal køres igen hvis input betingelser ændres en smule.

I dette projekt undersøges, hvordan Scientific Machine Learning (SciML) kan anvendes 
til at modellere og simulere plasmafysiske ligningssystemer. 

--- SciML har vist sig at kunne generalisere godt mellem gridpunkt af startbetingelser, så få træner+simulleringer kan generalisere rigtig godt ----

Projektet vil først fokusere på appalations versioner af HESEL-modellen,
til træning af state of the art (SOTA) SciML-metoder såsom Physics-Informed Neural Networks (PINN).
Herefter implimenteres og trænes disse modeller på det fulde HESEL-ligningssystem.

Målet er at undersøge, om SciML kan bidrage til mere fleksibel og effektiv modellering af 
plasmaets randdynamik.